\section{State of the Art}
\label{sec:intro_sota}

This section positions the thesis within the current landscape of modeling and simulation approaches relevant to heterogeneous hybrid system simulation.
It is a focused orientation that covers the main modeling paradigms, tool classes, and orchestration challenges that motivate the framework developed in this thesis.

%=================================================
\subsection{Equation-based Modeling and FMI-based Simulation Units}

Equation-based modeling describes systems through mathematical relations between physical quantities~\cite{Modelica_Fritzson}.
These relations may be algebraic equations, differential equations, or combinations of both.
Two modeling styles are especially relevant here.
In \emph{causal} modeling, the computational direction is defined through inputs and outputs.
In \emph{acausal} modeling, the model states relations between variables, and the simulation tool derives the computational causality.
Causal models are common for control, signal processing, and software-oriented subsystems.
Acausal equation-based models are useful for physical subsystems whose behavior follows conservation laws and constitutive relations.

Equation-based modeling is established for multi-domain physical systems.
It lets engineers describe subsystem behavior in a compact and domain-oriented form without prescribing the full numerical solution procedure.
This is useful for mechatronic systems because they combine mechanics, electronics, thermal effects, hydraulics, and control behavior.
A common equation-based formalism helps to represent these domains within one model structure and to connect them through physically meaningful interfaces.

The most prominent open language in this area is Modelica \cite{Modelica_Fritzson}.
Modelica is an object-oriented, equation-based language for modeling physical systems.
It supports hierarchical composition, reusable model libraries, and acausal connectors.
Tools such as Dymola, SimulationX, and OpenModelica~\cite{fritzson_openmodelica_2020} provide solvers, graphical editors, and libraries based on the Modelica language and the Modelica Standard Library.

For tool interoperability, the \ac{FMI} standard has become particularly important \cite{modelica_association_project_fmi_functional_2014}.
It defines how simulation units are packaged as \acp{FMU} and how an importing environment instantiates them, sets inputs and parameters, advances them through standardized calls, and reads their outputs at communication points.
This makes \ac{FMI} a widely used basis for co-simulation workflows across tools \cite{gomes_co-simulation_2019}.

%=================================================
\subsection{Domain-Specific Simulation Environments}

Equation-based modeling environments cover many technical domains, but some subsystem models require specialized modeling approaches and dedicated tools.
This is relevant for the application class considered in this thesis, where a technical device interacts with the human musculoskeletal system and where local structural effects can influence the system behavior.

One example is musculoskeletal simulation.
In applications such as prostheses and exoskeletons, the technical device is mechanically coupled to the human body \cite{eberle_biomechanical_2025}.
The resulting behavior depends on the device, the controller, and the biomechanics of the user.
OpenSim is a widely used open platform for musculoskeletal \ac{MBD}~\cite{OpenSim_Delp,seth_opensim_2018}.
It represents the human body as an articulated rigid-body system with joints, biomechanical constraints, and muscle-tendon actuation.
It also provides workflows such as forward dynamics and inverse dynamics.
OpenSim is therefore relevant for human-device interaction.
In co-simulation workflows, OpenSim models are usually integrated through the OpenSim \ac{API} or through backend-specific wrappers.

Finite element analysis is a second example.
Finite element models are used when spatially distributed quantities must be resolved explicitly.
This includes deformation, stress, strain, and local contact behavior.
Such quantities cannot usually be obtained from lumped rigid-body models alone.
The \ac{FEM} solves partial differential equations on a spatial discretization of the geometry and therefore provides a higher level of local detail \cite{knothe_finite_2017}.
This makes it suitable for compliant structures, structural loading, and contact phases in which local deformation is relevant.

Musculoskeletal and finite element models differ from classical equation-based models in how they expose their coupling interface.
Their coupling variables are often application-specific rather than described by a small standardized connector set.
They may include generalized coordinates, torques, external forces, contact quantities, marker trajectories, or reduced state variables.
As a result, these specialized model classes are often integrated through custom wrappers or tool-specific interfaces.
Direct \ac{FMU} export is less commonly available when solver control, state transfer, or application-specific coupling variables are required.

This creates a practical challenge for heterogeneous system simulation.
Domain-specific tools provide models that are needed for high-fidelity analysis, but they are not naturally embedded in \ac{FMU}-based co-simulation workflows.
A system simulation framework must therefore integrate both standardized \acp{FMU} and tool-specific simulation units under one common orchestration concept.

%=================================================
\subsection{Co-Simulation Frameworks and Orchestration}

Co-simulation enables the coupled simulation of subsystem models that originate from different tools.
Each subsystem keeps its own model representation and numerical solver.
A master algorithm coordinates the coupled simulation through data exchange at communication points.
This makes co-simulation a natural approach for heterogeneous system simulation \cite{gomes_co-simulation_2019}.

\ac{FMI} supports Model Exchange and Co-Simulation execution modes.
This thesis uses Co-Simulation \acp{FMU}, where each simulation unit keeps its own solver and is advanced through standardized calls at communication points.
\ac{FMI}~3.0 further adds scheduled execution and clock-based mechanisms for discrete-time behavior and event-oriented interaction \cite{modelica_association_project_fmi_functional_2024}.

Several frameworks build on \ac{FMI} to orchestrate coupled \ac{FMU} execution.
OMSimulator provides an \ac{FMI}-based composition and execution environment with scripting and graphical support~\cite{ochel_omsimulator_2019}.
INTO-CPS Maestro provides another orchestration environment for \ac{FMI}-based co-simulation and supports the specification and execution of coupled simulation scenarios~\cite{thule_maestro_2019}.
CoFMPy is a more recent Python-native framework for rapid prototyping of FMI-based Digital Twins, combining FMU co-simulation, algebraic-loop handling, datastream communication, and a Python FMU proxy for non-FMU components~\cite{friedrich_cofmpy_2025}.
These tools show that \ac{FMI}-based orchestration is established and that Python-native variants exist.

Within the Python ecosystem, lower-level libraries address the \ac{FMI} interface itself.
FMPy supports loading, inspecting, validating, and simulating \acp{FMU}~\cite{dassault_systemes__catia-systems_fmpy_2026}.
PyFMI provides a simulation-oriented Python interface for \emph{Model Exchange} and \emph{Co-Simulation} \acp{FMU}~\cite{andersson_pyfmi_2016}.
These libraries provide a mature basis for FMU access from Python.

The remaining difficulty is the broader heterogeneous workflow.
Direct interoperability is strongest when subsystem models are available as \acp{FMU}.
Models from tools that do not naturally export \acp{FMU}, such as musculoskeletal simulation environments or finite element solvers, usually require custom wrappers or tool-specific interfaces.
In addition, \ac{FMI} standardizes the simulation-unit interface, but it does not solve all orchestration problems at system level.
Execution order, algebraic loops, discrete events, and runtime switching between alternative subsystem models remain responsibilities of the importing framework.
These aspects define the gap addressed in this thesis.

%=================================================
\subsection{Hybrid Co-Simulation and Event Handling}

Hybrid systems combine continuous-time dynamics with discrete changes of state or behavior \cite{cellier_combined_1986,derler_modeling_2012}.
In \ac{CPS} and mechatronic systems, such behavior arises through sampled sensors, digital controllers, switching logic, mode changes, and mechanical contact.
Their simulation requires both continuous-time integration and consistent handling of discrete events.

In a monolithic simulation, one solver has access to the full system model.
It can evaluate all relevant variables and equations during each integration step.
This makes it possible to detect event conditions, localize event times, execute event logic, and continue from a consistent post-event state \cite{Modelica_Fritzson}.
For hybrid systems, this is a major advantage of a monolithic formulation.

In co-simulation, the master algorithm usually sees subsystem outputs only at communication points.
Each simulation unit advances internally between these points.
An event may therefore occur within a communication step without being immediately visible at system level.
Hybrid co-simulation requires additional orchestration mechanisms.
The importing framework must detect the event, localize its time, and restore or roll back subsystem states when necessary.
It must then execute the event handling logic, synchronize the affected simulation units, and resume the continuous simulation~\cite{broman_determinate_2013}.
The problem becomes more difficult when several events interact or when one event triggers further events in other subsystems.

The representation of time, event handling, and rollback capabilities have been identified as central issues for deterministic hybrid co-simulation \cite{StepRevision_Cremona,broman_requirements_2015,cremona_hybrid_2019}.
\ac{FMI}~3.0 adds concepts such as clocks, event mode, and scheduled execution \cite{modelica_association_project_fmi_functional_2024}.
These additions improve the representation of sampled and event-driven behavior.
The importing framework still has to implement the system-level logic for event detection, event localization, state consistency, and continuation after events.

Hybrid co-simulation is therefore central for the systems considered in this thesis.
Relevant examples include sensor sampling, controller action, mode changes, and mechanical contact.
A heterogeneous simulation framework must combine continuous integration with consistent event handling across coupled subsystem models.

%=================================================
\subsection{Multi-Model and Multi-Fidelity Simulation}

Many engineering subsystems can be represented by more than one model.
These models may differ in their physical assumptions, numerical complexity, computational cost, and resolved output quantities.
In multi-fidelity modeling, fidelity is a relative property of models that represent the same subsystem or input-output relation.
It describes the accuracy and physical detail provided for the quantities of interest, together with the computational cost required to obtain them~\cite{peherstorfer_survey_2018,fernandez-godino_review_2023}.
A multi-model simulation uses such alternatives within one workflow.

A low-dimensional or reduced-complexity model is often sufficient for global system behavior, controller design, or rapid parameter studies.
A detailed model becomes necessary when local effects dominate the quantity of interest.
Examples are deformation, stress, strain, and contact.
The appropriate model therefore depends on the simulation objective and on the current operating phase of the system~\cite{williams_switched-fidelity_2014,Choi2014}.

This trade-off appears both across the development process and within individual simulations.
Early design stages often rely on simple models because they are fast to build and simulate.
Later stages require more detailed models for analysis and verification.
Within one simulation scenario, different phases of motion may also require different levels of detail.

Region-based runtime model switching has been proposed as a strategy for reducing computational cost.
Choi et al. describe an approach in which a computationally expensive target model is replaced by lower-fidelity models outside an interest region while preserving a common external interface and switching according to fidelity-change conditions \cite{Choi2014,Choi2017}.
The present thesis is conceptually related to this idea, but focuses on heterogeneous subsystem models from different tools that share a common external interface.

Existing work on multi-fidelity simulation includes hierarchical model families, reduced-order models, surrogate models, and adaptive model selection \cite{peherstorfer_survey_2018,fernandez-godino_review_2023}.
For this thesis, the most relevant aspect is runtime switching between alternative subsystem models within one co-simulation workflow.
The main challenge is the consistent transition between models.
Such a workflow must preserve a common external interface, transfer relevant state information, define suitable switching criteria, and avoid unstable or excessively frequent mode changes \cite{williams_switched-fidelity_2014}.
In heterogeneous settings, the switch usually requires projection onto a shared reduced state rather than exact preservation of all internal model variables.
